\section{Reduced-form objects}

Throughout this section, fix a candidate $i$ and fix the opponents' profile $(\vminus,\aminus)$.
Let $\PhiStd(\cdot)$ denote the standard normal CDF.

\subsection{Opponents' strength aggregator}
\begin{definition}[Opponents' strength aggregator]\label{def:Sof}
For fixed $(\vminus,\aminus)$, define
\[
    \Sagg \equiv \Sof{\vminus}{\aminus} \;=\; \sum_{j\neq i} \ex{v_j + a_j}.
\]
\end{definition}

\begin{remark*}[Interpretation of $\Sagg$]
The scalar $\Sagg$ aggregates all information about candidate $i$'s opponents that is payoff-relevant for $i$'s primary nomination probability: holding $(v_i,a_i)$ fixed, $\Pr(w=i\mid\vprof,\aprof)$ depends on $(\vminus,\aminus)$ only through $\Sagg$.
\end{remark*}

\subsection{Nomination probabilities}
\begin{definition}[Reduced-form nomination probability]\label{def:pnomRF}
Fix $\Sagg>0$. For a candidate with valence $v\in\R$ who chooses action $a\in\{0,1\}$, define the probability of winning the primary as
\[
    \pnom{a}{v}{\Sagg} = \frac{\ex{v+a}}{\ex{v+a}+\Sagg}.
\]
We write
\[
    \pOne{v}{\Sagg} \equiv \pnom{1}{v}{\Sagg} = \pOneRF{v}{\Sagg},
    \qquad
    \pZero{v}{\Sagg} \equiv \pnom{0}{v}{\Sagg} = \pZeroRF{v}{\Sagg}.
\]
\end{definition}


The reduced-form map $\pnom{a}{v}{\Sagg}$ isolates how a candidate's own $(v,a)$ and the opponents' aggregate strength $\Sagg$ jointly determine the chance of winning the nomination.


\subsection{General-election win probabilities}
\begin{definition}[Reduced-form general-election win probability]\label{def:PgeRF}
Fix $\delta>0$. For a candidate with valence $v\in\R$ who chose action $a\in\{0,1\}$, define the probability of winning the general election as
\[
    \Pge{v}{\delta}{a} = \PhiStd(v-\delta a).
\]
\end{definition}


The parameter $\delta$ scales the electability penalty from choosing $a=1$ in the primary stage.


\subsection*{Expected utility and incentive to take action}
\begin{definition}[Reduced-form expected utility]\label{def:Urf}
Fix $(\Sagg,\delta)$ with $\Sagg>0$ and $\delta>0$. Define the expected payoff from choosing action $a\in\{0,1\}$ for a candidate with valence $v$ as
\[
    U_a(v;\Sagg,\delta) \equiv \pnom{a}{v}{\Sagg}\, \Pge{v}{\delta}{a}.
\]
\end{definition}


Since office requires winning both stages, $U_a(v;\Sagg,\delta)$ is the reduced-form probability of ultimately winning office when choosing action $a$.


\begin{definition}[Utility difference / incentive object]\label{def:Delta}
Fix $(\Sagg,\delta)$ with $\Sagg>0$ and $\delta>0$. Define
\[
    \Deltainc{v}{\Sagg}
    \equiv
    U_1(v;\Sagg,\delta) - U_0(v;\Sagg,\delta)
    =
    \pOne{v}{\Sagg}\PhiStd(v-\delta) - \pZero{v}{\Sagg}\PhiStd(v).
\]
\end{definition}


The sign of $\Deltainc{v}{\Sagg}$ determines whether taking action $a=1$ increases a candidate's chance of winning office, holding opponents fixed.


\subsection{Proportional gain and cost objects}
\begin{definition}[Proportional nomination gain]\label{def:Gbasic}
Fix $\Sagg>0$. Define the proportional nomination gain from choosing $a=1$ rather than $a=0$ as
\[
    \Gfun{v}{\Sagg} \equiv \frac{\pOne{v}{\Sagg}}{\pZero{v}{\Sagg}}.
\]
\end{definition}


The object $\Gfun{v}{\Sagg}$ isolates the primary-stage benefit of choosing $a=1$ in proportional terms.


\begin{definition}[Proportional electability cost]\label{def:Cbasic}
Fix $\delta>0$. Define the proportional electability cost of choosing $a=1$ rather than $a=0$ as
\[
    \Cfun{v}{\delta} \equiv \frac{\PhiStd(v)}{\PhiStd(v-\delta)}.
\]
\end{definition}


The object $\Cfun{v}{\delta}$ isolates the general-election penalty of choosing $a=1$ in proportional terms.
